% ============================================================
\section{Application Domains}
% ============================================================

\subsection{Sleep position monitoring}

The monitoring of spinal posture during sleep represents the primary disruptive
application for direct 3D shape reconstruction, and the one in which the technology's
advantages over current alternatives are most pronounced.

\subsubsection{Methodological limitations of indirect mattress sensing}

Consumer systems and research mattress sensors infer sleeping posture from pressure
distributions beneath the body surface, estimating whether a person lies on their left
or right side.
This approach characterises the mechanics of the mattress surface, not of the spine
itself.
Mattress pressure reflects the distribution of body mass across the support surface;
it does not measure vertebral curvature, inter-segment torsion, or lumbosacral
lordosis angle.
In orthopaedic terms, these mattress-derived classifications are methodological
workarounds: the variables of clinical interest are in the spine, and the sensor is
in the mattress.
Optical cameras placed above the bed face equivalent limitations --- they estimate
trunk pose from external silhouette data, are sensitive to bedding occlusion, and
cannot resolve inter-vertebral angles at clinical precision.

\subsubsection{Direct spinal measurement during sleep}

Placed in a compression shirt overnight, FlexTail measures five spinal metrics
continuously, none of which is accessible to mattress or camera systems:

\begin{enumerate}
	\item Lumbar curvature (L-spine angle in the sagittal plane).
	\item Lateral flexion of the lumbar and lower thoracic spine.
	\item Localised axial torsion (inter-vertebral rotation about the longitudinal axis).
	\item Upper thoracic curvature (shoulder-girdle level).
	\item Lumbosacral lordosis and thoracic kyphosis angles.
\end{enumerate}

\subsubsection{Drift advantage in the sleep context}

Yaw drift --- the primary dynamic limitation of IMU-based systems in high-speed
activity --- is substantially reduced in the sleep context.
Postural transitions during sleep occur over minutes rather than seconds.
The slow rate of change allows the \emph{flexlib} software to identify drift
trajectories precisely and correct them post-hoc, producing accurate curvature
records even over multi-hour recordings without requiring magnetometer input.

\subsubsection{Automated cloud-based data workflow}

Longitudinal sleep studies traditionally require manual data collection: researchers
visit participants periodically to extract files from recording devices.
The FlexTail system automates this workflow.
The participant activates the sensor, pairs it with the smartphone application, and
begins sleep.
Data are transmitted to the cloud infrastructure during or after the recording
session, eliminating manual file transfers and the associated transcription errors.

\subsection{Occupational health: construction ergonomics}

The clearest published demonstration of FlexTail's advantage in uncontrolled
environments comes from Sawicki, D{\"u}king et al.\ (2026) at TU Braunschweig
\cite{sawicki2026}.
The study compared musculoskeletal loading in two concrete construction processes:
conventional manual casting and a novel collaborative \ac{SC3DP}
process in which workers guide a robotic arm rather than handle materials directly.

The study's measurement stack illustrates precisely why a dedicated spinal sensor is
necessary even when full optical infrastructure is available.
Sawicki et al.\ deployed seven Vicon Valkyrie VK16 cameras covering the entire working
volume, with retroreflective targets fixed to each worker's helmet.
This system successfully tracked head position throughout the production runs and
yielded the covered-distance figures reported in the results.
It provided no information about the spine.
Helmet-marker trajectory records where a worker walked; it does not record whether
the lumbar spine was flexed, how far, or for how long.
Recovering vertebral curvature from optical data would require dense marker arrays
mounted directly over each spinal segment --- impractical on workers handling
reinforcement bars and concrete hose in a live fabrication environment.
FlexTail, worn inside the standard shirt, delivered continuous spinal posture data
under exactly those conditions, yielding the 60\,\% reduction in flexed-posture time
that is the study's principal ergonomic finding.

A second methodological advantage is relevant here: because the Sensorshirt is
indistinguishable from ordinary workwear, participants are unaware of ongoing kinematic
recording during normal work.
This eliminates the Hawthorne effect --- the well-documented modification of behaviour
that occurs when participants know they are being observed during short ergonomic
assessments --- and produces kinematic data representative of habitual practice.

Results showed that \acs{SC3DP} reduced:

\begin{itemize}
	\item Time in flexed spinal postures by \textbf{60\,\%}.
	\item Total carried weight by \textbf{44\,\%}.
	\item Distance walked per task by \textbf{37\,\%}.
	\item Perceived exertion (Borg and NASA-TLX scores) by \textbf{63\,\%}.
\end{itemize}

Objective physiological markers (heart rate, blood lactate) were unchanged between
conditions, indicating that SC3DP substituted physical for cognitive demands rather
than reducing total strain.
Task-level productivity was approximately three times higher in the robotic condition.

The study exemplifies a methodological pattern likely to recur across construction,
logistics, mining, and process industries: optical infrastructure can be present and
still leave spinal ergonomics unmeasured, because marker placement on a moving torso
is incompatible with physical work.
Wearable shape reconstruction fills that specific gap without displacing other
instrumentation.

\subsection{Clinical rehabilitation and post-surgical follow-up}

Standard post-surgical assessment relies on static in-clinic \acs{ROM} tests and
patient-reported questionnaires.
The sensitivity of both methods is limited: ROM tests capture behaviour during a brief,
supervised session; questionnaires are subject to recall bias, social desirability
effects, and response shift \cite{natarajan2023}.

Continuous ambulatory recording with FlexTail enables a different approach.
The raw curvature stream contains all information needed to extract standardised ROM
metrics: maximum segmental flexion, extension, lateral flexion, and torsion angles can
be computed automatically from any movement episode that elicits them, whether during
a formal test or spontaneously during daily life.
The software also computes a calibrated posture score --- a composite metric weighted
by the duration and severity of postural deviations from a neutral reference --- that
tracks functional recovery over weeks without requiring repeated clinical visits.

The Lindenhofgruppe (Switzerland), under PD Dr.\ med.\ Emin Aghayev, is currently
using the technology to quantify spine mobility trajectories in free-living patients
following spinal surgery \cite{marx2023}.
This collaboration is ongoing; no peer-reviewed publication from the free-living phase
is available at the time of writing.

\subsection{Sport biomechanics and fitness tracking}

In resistance training and sport, the clinically meaningful quantity is not whether
a movement is performed but \emph{how}.
The FlexTail sensor quantifies execution quality in exercises where the spine is the
primary measurement site.

\subsubsection{Squat: pelvic tilt detection and longitudinal progression}

During a barbell back squat, lumbar neutral must be maintained through the descent.
Loss of lumbar lordosis (``butt wink'') at depth increases compressive and shear loads
on the posterior disc and facet joints.
The FlexTail base-IMU records vertical displacement as the proxy for squat depth;
the strain-gauge array records lumbar curvature simultaneously.
The software reports the precise percentage of descent at which neutral is lost for
each individual repetition, providing trainers and physiotherapists with an objective
time course of lumbar control rather than a binary pass/fail judgement.

Longitudinal tracking quantifies physiological adaptation directly.
In a pilot observation, a participant who initially lost lumbar neutrality at 50\,\% of
squat depth maintained neutral to 90\,\% of depth following an eight-week training
intervention --- an objective kinematic outcome that could not have been derived from
self-report or visual inspection alone.

The same principle applies to push-up lumbar sag (where the sensor detects progressive
extension loss across a set) and overhead-press spinal extension under fatigue.
Bench press, performed in strict supine position, is the principal exercise where
the dorsal placement of the sensor limits measurement utility.

\subsubsection{Movement archetypes}

The five movement archetypes framework --- hip-hinge, squat, push, pull, carry ---
maps naturally to the continuous curvature and torsion signal from the sensor array.
Deadlift and Romanian deadlift generate characteristic sagittal flexion-extension
curves; overhead press generates thoracolumbar extension; loaded rows generate
asymmetric torsion signals.
Each archetype produces a kinematically distinct fingerprint that can be used to
train activity-classification models analogous to those validated by Walkling et al.\
for \acs{ADL} recognition \cite{walkling2025}.

\subsubsection{High-speed sport: golf and equestrian}

Preliminary observations in golf demonstrated that spinal posture errors during the
backswing and follow-through (e.g., reverse spine angle) are distinguishable from
correct mechanics by inspection of the raw curvature trace.
In an experimental setup with amateur golfers, the technology differentiated incorrect
from correct swings by visual inspection of generated data traces; automated
machine-learning classification of such errors is under development.

Equestrian sport --- particularly dressage --- presents a complementary case where
spinal posture is not incidental but constitutively scored.
FEI judging criteria for collected movements and piaffe explicitly penalise rider
collapse in the lumbar spine, lateral trunk lean, and behind-the-vertical seat
position, all of which produce distinct curvature and torsion signatures in the
sensor array.
Unlike the ballistic constraints of a golf swing, dressage posture is sustained
across movements lasting tens of seconds, placing it well within the validated
velocity envelope of the current sensor stack.
Continuous recording across a training session can therefore quantify postural
compliance with competition standards, identify fatigue-related collapse in
thoracolumbar extension, and provide objective feedback to replace the
inherently intermittent input of a trainer observing from the arena.

Automated classification of ballistic sport errors --- relevant to the golf case ---
requires handling the high angular
accelerations that produce significant short-term \acs{IMU} transients.
Sport-specific filter tuning to address this is an active area of development, and the
current software stack is not validated for movements faster than approximately
60\,\si{\degree\per\second}; this constraint does not apply to equestrian posture
monitoring, where trunk angular velocities remain well below this threshold.

\subsection{Occupational health: nursing and care work}

Musculoskeletal disorders account for a disproportionate share of sick-leave days in
nursing and care professions \cite{sun2023}.
The specific risk factors --- sustained lumbar flexion while attending to patients in
low hospital beds, asymmetric load transfer during manual patient handling, and
repetitive trunk torsion --- are well characterised in the literature but poorly
quantified in individual workers' actual shifts.

Conventional short-term ergonomic workplace observations suffer from the Hawthorne
effect: workers modify their movement behaviour when they are aware of being observed,
systematically underrepresenting habitual postural load.
Longitudinal wearable monitoring with the Sensorshirt eliminates this confound.
Workers wear the shirt throughout normal shifts without modifying routine; the
recording is continuous, invisible under standard uniforms, and extends over weeks.

Analysis of these recordings in FlexLab Studio consistently reveals sharp
discrepancies between workers' subjective perception of their postural load and the
objectively measured biomechanical exposure: participants frequently underestimate the
duration of sustained hyperkyphosis and the torsional load during patient transfers.

Pilot programmes integrating the Sensorshirt into occupational health management in
collaboration with major health insurance funds have demonstrated a measurable
intervention effect.
Vibration feedback from the onboard motor, triggered when cumulative time in
hyperkyphosis or high torsion exceeds configurable thresholds, increased the proportion
of working time spent in an upright spinal posture from approximately 10\,\% to
15\,\% in one pilot cohort.
This improvement was associated with a reduction in self-reported acute low-back pain
among participants.
These results are from internal pilot data; peer-reviewed validation is in progress.

\subsection{Space medicine}

In microgravity, the absence of axial compressive loading allows spinal discs to
absorb fluid, elongating the spine by several centimetres \cite{belavy2016}.
Simultaneously, paraspinal musculature atrophies without gravitational loading stimulus.
This combination predisposes returning astronauts to acute disc herniation during
re-exposure to normal gravity: Belavy et al.\ (2016) document elevated herniation risk
in the post-flight period compared with pre-flight baseline \cite{belavy2016}.

Mapping how vertebral curvature, inter-segment flexibility, and muscle-supported
posture change across simulated microgravity (such as head-down tilt bed rest) requires
a sensor that can be applied longitudinally without radiation exposure and without a
motion-capture laboratory.

MinkTec is collaborating with the UZH Space Hub and the \ac{ESA} on this question.
The technology is being evaluated in the ThromBoShift bed-rest study, which uses
head-down tilt as a controlled analogue of spaceflight-induced spinal unloading.
Because the compression shirt cannot be worn during prolonged bed rest (risk of
skin-surface displacement under continuous dorsal loading), the protocol uses
medical-grade adhesive patches to fix the sensor strip directly to the skin,
preserving full mechanical coupling.
No peer-reviewed data from this collaboration are available at the time of writing.
